\chapter{Danh mục tài liệu tham khảo}\label{ch:11}

\section*{Tài liệu tiếng Việt}
\begin{list}{}{\setlength{\leftmargin}{1.5em}\setlength{\itemindent}{-1.5em}\setlength{\itemsep}{2pt}\small}
\item Báo Đầu tư Chứng khoán. (2026). HNX nhập HOSE, UPCoM còn doanh nghiệp nào đang chia cổ tức tốt với P/E rẻ? [Bài báo ngày 06/10/2026].
\item Bộ Tài chính. (2020). Thông tư số 120/2020/TT-BTC ngày 31/12/2020 quy định giao dịch cổ phiếu niêm yết, đăng ký giao dịch và chứng chỉ quỹ, trái phiếu doanh nghiệp, chứng quyền có bảo đảm niêm yết trên hệ thống giao dịch chứng khoán.
\item Bộ Tài chính. (2023). Văn bản hợp nhất số 19/VBHN-BTC ngày 25/12/2023 hợp nhất Thông tư quy định lộ trình sắp xếp lại thị trường giao dịch cổ phiếu, thị trường giao dịch trái phiếu, thị trường giao dịch chứng khoán phái sinh và thị trường giao dịch các loại chứng khoán khác (hợp nhất Thông tư số 57/2021/TT-BTC ngày 12/7/2021 và Thông tư số 69/2023/TT-BTC ngày 15/11/2023).
\item Bộ Tài chính. (2025). Thông tư số 139/2025/TT-BTC ngày 30/12/2025 sửa đổi, bổ sung một số điều của Thông tư số 57/2021/TT-BTC ngày 12 tháng 7 năm 2021 của Bộ trưởng Bộ Tài chính quy định lộ trình sắp xếp lại thị trường giao dịch cổ phiếu, thị trường giao dịch trái phiếu, thị trường giao dịch chứng khoán phái sinh và thị trường giao dịch các loại chứng khoán khác được sửa đổi, bổ sung bởi Thông tư số 69/2023/TT-BTC ngày 15 tháng 11 năm 2023 của Bộ trưởng Bộ Tài chính.
\item Hội đồng nhân dân thành phố Hà Nội. (2026). Nghị quyết số 57/2026/NQ-HĐND ngày 15/6/2026 quy định một số chính sách phát triển nguồn nhân lực chất lượng cao của thành phố Hà Nội.
\item Nguyễn, T. B., \& Nguyễn, V. T. (2026b). Tiền gửi không kỳ hạn, hiệu quả hoạt động và ổn định tài chính ngân hàng: Bằng chứng từ mô hình ngưỡng tại Việt Nam. \emph{Tạp chí Kinh tế - Luật và Ngân hàng, 28}(9), 52–64. \url{https://doi.org/10.59276/3030-4199/jelb.bav.895}
\item Quốc hội. (2025). Nghị quyết số 222/2025/QH15 ngày 27/6/2025 về Trung tâm tài chính quốc tế tại Việt Nam.
\item Sở Giao dịch Chứng khoán Hà Nội. (2016). Quyết định số 529/QĐ-SGDHN ngày 25/8/2016 về việc ban hành Quy chế giao dịch chứng khoán niêm yết tại Sở Giao dịch Chứng khoán Hà Nội.
\item Sở Giao dịch Chứng khoán Thành phố Hồ Chí Minh. (2016). Quyết định số 341/QĐ-SGDHCM ngày 22/8/2016 ban hành Quy chế giao dịch chứng khoán.
\item Sở Giao dịch Chứng khoán Thành phố Hồ Chí Minh. (2025). Quy định giao dịch đăng website HOSE [Kèm theo Tờ trình số 51/TTr-HTGD ngày 26/4/2025].
\item Sở Giao dịch Chứng khoán Thành phố Hồ Chí Minh. (2026). Quyết định số 737/QĐ-SGDHCM ngày 29/9/2026 ban hành Quy tắc chỉ số VNINDEX phiên bản 1.0. \url{https://stream.vnx.vn/VNX//Article/2026100813094820920260929_QD737_SGDHCM_QTCS-VNINDEX.pdf}
\item Sở Giao dịch Chứng khoán Việt Nam. (2026). Sắp xếp lại thị trường giao dịch cổ phiếu niêm yết theo Thông tư số 139/2025/TT-BTC ngày 30/12/2025 của Bộ Tài chính [Thông báo ngày 30/9/2026]. \url{https://vnx.vn/vi/thong-tin-tu-sgdck/VNX_QLHD/10005421}
\item Thời báo Tài chính Việt Nam. (2026). Chuyển toàn bộ cổ phiếu HNX sang HOSE trong năm 2026, nghiên cứu giao dịch xuyên trưa và nới biên độ. \url{https://thoibaotaichinhvietnam.vn/chuyen-toan-bo-co-phieu-hnx-sang-hose-trong-nam-2026-nghien-cuu-giao-dich-xuyen-trua-va-noi-bien-do-204943.html}
\item Tuổi Trẻ Online. (2026a). HOSE dự kiến chia cổ phiếu thành hai bảng Prime Board và Standard Board. \url{https://tuoitre.vn/hose-du-kien-chia-co-phieu-thanh-hai-bang-prime-board-va-standard-board-100261002092239776.htm}
\item Tuổi Trẻ Online. (2026b). Toàn bộ cổ phiếu trên HNX sẽ chuyển sang HOSE từ 28-12 [Bài báo ngày 30/9/2026]. \url{https://tuoitre.vn/toan-bo-co-phieu-tren-hnx-se-chuyen-sang-hose-tu-28-12-100260930201353937.htm}
\item Ủy ban nhân dân thành phố Hà Nội. (2026). Quyết định số 1193/QĐ-UBND ngày 18/3/2026 về việc công bố Danh mục các bài toán lớn của thành phố Hà Nội giai đoạn 2026–2030 (đợt 1).
\end{list}
\section*{Tài liệu tiếng Anh}
\begin{list}{}{\setlength{\leftmargin}{1.5em}\setlength{\itemindent}{-1.5em}\setlength{\itemsep}{2pt}\small}
\item Amihud, Y. (2002). Illiquidity and stock returns: Cross-section and time-series effects. \emph{Journal of Financial Markets, 5}(1), 31–56.
\item Amihud, Y., Mendelson, H., \& Lauterbach, B. (1997). Market microstructure and securities values: Evidence from the Tel Aviv Stock Exchange. \emph{Journal of Financial Economics, 45}(3), 365–390. \url{https://doi.org/10.1016/S0304-405X(97)00021-4}
\item Bae, K.-H., Chan, K., \& Ng, A. (2004). Investibility and return volatility. \emph{Journal of Financial Economics, 71}(2), 239–263.
\item Bai, J., \& Perron, P. (1998). Estimating and testing linear models with multiple structural changes. \emph{Econometrica, 66}(1), 47–78.
\item Bekaert, G., Harvey, C. R., \& Lundblad, C. (2007). Liquidity and expected returns: Lessons from emerging markets. \emph{The Review of Financial Studies, 20}(6), 1783–1831.
\item Bessler, W., Opfer, H., \& Wolff, D. (2017). Multi-asset portfolio optimization and out-of-sample performance: An evaluation of Black-Litterman, mean-variance, and naïve diversification approaches. \emph{The European Journal of Finance, 23}(1), 1–30.
\item Black, F., \& Litterman, R. (1992). Global portfolio optimization. \emph{Financial Analysts Journal, 48}(5), 28–43.
\item Callaway, B., \& Sant’Anna, P. H. C. (2021). Difference-in-differences with multiple time periods. \emph{Journal of Econometrics, 225}(2), 200–230.
\item Cetorelli, N., \& Peristiani, S. (2013). Prestigious stock exchanges: A network analysis of international financial centers. \emph{Journal of Banking \& Finance, 37}(5), 1543–1551.
\item Chen, K., \& Chen, G. (2015). The rise of international financial centers in mainland China. \emph{Cities, 47}, 10–22.
\item Christie, W. G., \& Huang, R. D. (1994). Market structures and liquidity: A transactions data study of exchange listings. \emph{Journal of Financial Intermediation, 3}(3), 300–326.
\item Čihák, M., Demirgüç-Kunt, A., Feyen, E., \& Levine, R. (2012). \emph{Benchmarking financial systems around the world} (Policy Research Working Paper No. 6175). World Bank.
\item Conley, T. G., \& Taber, C. R. (2011). Inference with \enquote{difference in differences} with a small number of policy changes. \emph{The Review of Economics and Statistics, 93}(1), 113–125.
\item Corwin, S. A., \& Schultz, P. (2012). A simple way to estimate bid-ask spreads from daily high and low prices. \emph{The Journal of Finance, 67}(2), 719–760. \url{https://doi.org/10.1111/j.1540-6261.2012.01729.x}
\item Coval, J. D., \& Moskowitz, T. J. (1999). Home bias at home: Local equity preference in domestic portfolios. \emph{The Journal of Finance, 54}(6), 2045–2073. \url{https://doi.org/10.1111/0022-1082.00181}
\item Dao, V. H., Nguyen, T. M. H., \& Vu, T. D. (2021). The impact of COVID-19 on stock market returns in Vietnam. \emph{Journal of Risk and Financial Management, 14}(9), Article 441. \url{https://doi.org/10.3390/jrfm14090441}
\item Daw, J. R., \& Hatfield, L. A. (2018). Matching and regression to the mean in difference-in-differences analysis. \emph{Health Services Research, 53}(6), 4138–4156. \url{https://doi.org/10.1111/1475-6773.12993}
\item DeMiguel, V., Garlappi, L., \& Uppal, R. (2009). Optimal versus naive diversification: How inefficient is the 1/N portfolio strategy? \emph{The Review of Financial Studies, 22}(5), 1915–1953.
\item Eichengreen, B., \& Shah, N. (2020). The correlates of international financial-center status. \emph{Review of International Economics, 28}(1), 62–81. \url{https://doi.org/10.1111/roie.12441}
\item Elliott, D. J. (2011). \emph{Building a global financial center in Shanghai: Observations from other centers}. Brookings Institution.
\item Fong, K. Y. L., Holden, C. W., \& Trzcinka, C. A. (2017). What are the best liquidity proxies for global research? \emph{Review of Finance, 21}(4), 1355–1401.
\item FTSE Russell. (2026a, September). \emph{FTSE Global Equity Index Series ground rules} (Version 14.4). LSEG.
\item FTSE Russell. (2026b, August). \emph{Reclassification of Vietnam from Frontier to Secondary Emerging market status: Frequently asked questions} (Version 1.3). LSEG. \url{https://www.lseg.com/content/dam/ftse-russell/en_us/documents/policy-documents/ftse-faq-document-vietnam-reclassification.pdf}
\item Garlappi, L., Uppal, R., \& Wang, T. (2007). Portfolio selection with parameter and model uncertainty: A multi-prior approach. \emph{The Review of Financial Studies, 20}(1), 41–81.
\item Gehrig, T. (2000). Cities and the geography of financial centers. In J.-M. Huriot \& J.-F. Thisse (Eds.), \emph{Economics of cities: Theoretical perspectives} (pp. 415–445). Cambridge University Press.
\item Goutte, S., Le, H.-V., Liu, F., \& von Mettenheim, H.-J. (2025). MCDA strategies for portfolio optimization: A case study on Vietnamese stock market dynamics. \emph{Annals of Operations Research, 353}(1), 321–351.
\item Griffin, J. M., Kelly, P. J., \& Nardari, F. (2010). Do market efficiency measures yield correct inferences? A comparison of developed and emerging markets. \emph{The Review of Financial Studies, 23}(8), 3225–3277.
\item Hamilton, J. D. (1989). A new approach to the economic analysis of nonstationary time series and the business cycle. \emph{Econometrica, 57}(2), 357–384.
\item Hisada, T. (2018). The effect of investor sentiment toward an exchange merger on liquidity. \emph{International Journal of Economics and Financial Issues, 8}(2), 315–318.
\item Hoang, V.-M. D., \& Luu, D. T. (2024). \emph{Modelling volatility dynamics and identifying regime switching in a frontier stock market: The case of Vietnam} (SSRN Working Paper No. 4867203; bản sửa ngày 17/9/2025). SSRN.
\item Holm, S. (1979). A simple sequentially rejective multiple test procedure. \emph{Scandinavian Journal of Statistics, 6}(2), 65–70.
\item Hou, K., \& Moskowitz, T. J. (2005). Market frictions, price delay, and the cross-section of expected returns. \emph{The Review of Financial Studies, 18}(3), 981–1020.
\item Huang, R. D., \& Stoll, H. R. (1996). Dealer versus auction markets: A paired comparison of execution costs on NASDAQ and the NYSE. \emph{Journal of Financial Economics, 41}(3), 313–357.
\item International Monetary Fund. (2023). \emph{Vietnam: 2023 Article IV consultation} (IMF Country Report No. 23/338). \url{https://www.elibrary.imf.org/view/journals/002/2023/338/article-A001-en.xml}
\item Kan, R., \& Zhou, G. (2007). Optimal portfolio choice with parameter uncertainty. \emph{Journal of Financial and Quantitative Analysis, 42}(3), 621–656.
\item Kaufmann, D., Kraay, A., \& Mastruzzi, M. (2011). The Worldwide Governance Indicators: Methodology and analytical issues. \emph{Hague Journal on the Rule of Law, 3}(2), 220–246.
\item Kindleberger, C. P. (1974). \emph{The formation of financial centers: A study in comparative economic history} (Princeton Studies in International Finance No. 36). Princeton University.
\item Kritzman, M., Page, S., \& Turkington, D. (2012). Regime shifts: Implications for dynamic strategies. \emph{Financial Analysts Journal, 68}(3), 22–39.
\item Lai, K. P. Y. (2012). Differentiated markets: Shanghai, Beijing and Hong Kong in China’s financial centre network. \emph{Urban Studies, 49}(6), 1275–1296. \url{https://doi.org/10.1177/0042098011408143}
\item Lang, M., Lins, K. V., \& Maffett, M. (2012). Transparency, liquidity, and valuation: International evidence on when transparency matters most. \emph{Journal of Accounting Research, 50}(3), 729–774.
\item Ledoit, O., \& Wolf, M. (2004). A well-conditioned estimator for large-dimensional covariance matrices. \emph{Journal of Multivariate Analysis, 88}(2), 365–411.
\item Ledoit, O., \& Wolf, M. (2008). Robust performance hypothesis testing with the Sharpe ratio. \emph{Journal of Empirical Finance, 15}(5), 850–859.
\item Lesmond, D. A. (2005). Liquidity of emerging markets. \emph{Journal of Financial Economics, 77}(2), 411–452.
\item Lesmond, D. A., Ogden, J. P., \& Trzcinka, C. A. (1999). A new estimate of transaction costs. \emph{The Review of Financial Studies, 12}(5), 1113–1141.
\item Lo, A. W., \& MacKinlay, A. C. (1988). Stock market prices do not follow random walks: Evidence from a simple specification test. \emph{The Review of Financial Studies, 1}(1), 41–66.
\item Loughran, T., \& Schultz, P. (2005). Liquidity: Urban versus rural firms. \emph{Journal of Financial Economics, 78}(2), 341–374.
\item Markowitz, H. (1952). Portfolio selection. \emph{The Journal of Finance, 7}(1), 77–91.
\item Michaud, R. O. (1989). The Markowitz optimization enigma: Is 'optimized' optimal? \emph{Financial Analysts Journal, 45}(1), 31–42.
\item Moosa, I., Li, L., \& Jiang, R. (2016). Determinants of the status of an international financial centre. \emph{The World Economy, 39}(12), 2074–2096.
\item Naeem, M. M. F. (2023). \emph{What makes an International Financial Centre (IFC) competitive? An empirical study of the determinants responsible for the competitiveness of an IFC} [Luận án Tiến sĩ Quản trị kinh doanh (DBA), Durham University]. Durham E-Theses. \url{https://etheses.durham.ac.uk/id/eprint/15102/}
\item Nguyen, N., Nguyen, T., Tran, T., \& Mai, A. (2020). Shrinkage model selection for portfolio optimization on Vietnam stock market. \emph{The Journal of Asian Finance, Economics and Business, 7}(9), 135–145.
\item Nguyễn, T. B., \& Nguyễn, V. T. (2026a). Black-Litterman portfolio optimization using regime switching CAPM and ABC-MCMC: Empirical evidence from the Vietnamese stock market period 2019–2025. \emph{Tạp chí Nghiên cứu Chính sách và Phát triển, 3}(1), 81–99. \url{https://doi.org/10.63640/3030-4091/jpd.apd.194}
\item Nielsson, U. (2009). Stock exchange merger and liquidity: The case of Euronext. \emph{Journal of Financial Markets, 12}(2), 229–267.
\item OECD. (2008). \emph{Handbook on constructing composite indicators: Methodology and user guide}. OECD Publishing.
\item Pflug, G. C., Pichler, A., \& Wozabal, D. (2012). The 1/N investment strategy is optimal under high model ambiguity. \emph{Journal of Banking \& Finance, 36}(2), 410–417.
\item Politis, D. N., \& Romano, J. P. (1994). The stationary bootstrap. \emph{Journal of the American Statistical Association, 89}(428), 1303–1313.
\item Rambachan, A., \& Roth, J. (2023). A more credible approach to parallel trends. \emph{The Review of Economic Studies, 90}(5), 2555–2591.
\item Ta, B. Q., \& Vuong, T. (2020). The Black-Litterman model for portfolio optimization on Vietnam stock market. \emph{International Journal of Uncertainty, Fuzziness and Knowledge-Based Systems, 28}(Supp01), 99–111.
\item Vo, X. V., \& Phan, D. B. A. (2017). Further evidence on the herd behavior in Vietnam stock market. \emph{Journal of Behavioral and Experimental Finance, 13}, 33–41.
\item Vu, T. L., Luong, N. H., \& Nguyen, P. N. (2024). Stock crash and how it spreads over the Vietnamese stock market. \emph{VNU Journal of Economics and Business, 4}(3), 107–116. \url{https://doi.org/10.57110/vnu-jeb.v4i3.339}
\item Wang, D. T., Zhao, S. X., \& Wang, D. (2007). \enquote{Information hinterland}: A base for financial centre development. The case of Beijing versus Shanghai in China. \emph{Tijdschrift voor Economische en Sociale Geografie, 98}(1), 102–120.
\item Wardle, M., \& Mainelli, M. (2026, September). \emph{The Global Financial Centres Index 40}. Z/Yen Group \& China Development Institute. \url{https://www.longfinance.net/media/documents/GFCI_40_Report_2026.09.16_v1.2.pdf}
\item Wójcik, D. (2007). Geography and the future of stock exchanges: Between real and virtual space. \emph{Growth and Change, 38}(2), 200–223.
\item Yao, J.-M. (2012). \emph{Compétitivité et attractivité des places financières internationales: Application à l’industrie des fonds d’investissement au Luxembourg} [Luận án tiến sĩ Kinh tế học, Université de Strasbourg].
\end{list}

\vspace{2\baselineskip}
\begin{flushright}
\begin{minipage}{0.45\textwidth}\centering
Hà Nội, ngày \dots\ tháng 10 năm 2026\\
\textsc{Người dự tuyển}\\
(Ký và ghi rõ họ tên)\\[3\baselineskip]
Nguyễn Văn Trung
\end{minipage}
\end{flushright}
